% !TeX root = ../main.tex

\section{Краткое функциональное и математическое описание устройства}

Проектируемое устройство предназначено для выделения, цифровой фильтрации и децимации радиосигнала и состоит из квадратурного демодулятора и многоскоростной системы фильтрации.

\subsection{Квадратурная демодуляция}
Входной аналоговый сигнал оцифровывается аналого-цифровым преобразователем (АЦП) на высокой частоте дискретизации $f_s$. Для переноса спектра сигнала на околонулевые частоты (видеочастоту) и разделения его на синфазную ($I$) и квадратурную ($Q$) составляющие, оцифрованный сигнал перемножается с ортогональными колебаниями, генерируемыми цифровым гетеродином (ЦГ). 
Математически формирование квадратур описывается следующим образом:
\begin{equation}
    I(n) = x(n) \cdot \cos(2\pi f_G n / f_s)
\end{equation}
\begin{equation}
    Q(n) = x(n) \cdot \sin(2\pi f_G n / f_s)
\end{equation}
где $x(n)$ — отсчеты входного сигнала, $f_G$ — центральная частота цифрового гетеродина.

\subsection{Интегрально-гребенчатый фильтр (CIC-фильтр)}
Для эффективного подавления внеполосных спектральных составляющих и понижения частоты дискретизации (децимации) в $R$ раз применяется каскадный интегрально-гребенчатый фильтр (фильтр Хогенауэра). Его вычислительное преимущество заключается в том, что он не требует выполнения операций умножения, а реализуется исключительно на сумматорах и линиях задержки.

В рамках работы используется фильтр 8-го порядка ($N=8$) с коэффициентом децимации $R=10$. Фильтр состоит из $N$ каскадов интеграторов (работающих на частоте $f_s$), блока дециматора и $N$ каскадов гребенчатых фильтров (работающих на пониженной частоте).

Разностное уравнение одиночного цифрового фильтра-интегратора имеет вид:
\begin{equation}
    y(n) = y(n-1) + x(n)
\end{equation}
а его передаточная функция в $z$-области:
\begin{equation}
    H_I(z) = \frac{1}{1 - z^{-1}}
\end{equation}
где $z^{-1} = \exp(-j\omega)$ — оператор единичной задержки.

Общая передаточная функция CIC-фильтра $N$-го порядка описывается выражением:
\begin{equation}
    H_{CIC}(z) = \left( \frac{1 - z^{-D}}{1 - z^{-1}} \right)^N
\end{equation}
где $D$ — коэффициент дифференциальной задержки[cite: 2].

\subsection{КИХ-фильтр компенсатор}
Недостатком CIC-фильтра является спад амплитудно-частотной характеристики (АЧХ) в рабочей полосе пропускания, что приводит к амплитудным искажениям полезного сигнала. Для компенсации этого спада, а также для формирования узкой переходной зоны и требуемого уровня подавления вне полосы, используется дополнительный фильтр с конечной импульсной характеристикой (КИХ-корректор). 
КИХ-корректор устанавливается на выходе гребенчатых фильтров и работает на пониженной частоте дискретизации, а его АЧХ формируется таким образом, чтобы компенсировать завал АЧХ фильтра Хогенауэра. 

В пункте 7 работы каскад из CIC-фильтра и КИХ-корректора заменяется на единый эквивалентный КИХ-фильтр дециматор с аналогичными параметрами подавления, который осуществляет всю фильтрацию и децимацию в один шаг.